% ========================================================
% BÀI TẬP TOÁN 5 - KẾT NỐI TRI THỨC VỚI CUỘC SỐNG
% Bài 62: Luyện tập chung
% Định dạng: BTTeX Part III - Trắc nghiệm điền số ngắn
% ========================================================

% Câu 1
\begin{ex}
  Tính: $3\text{ giờ } 25\text{ phút} + 2\text{ giờ } 45\text{ phút}$. Đổi ra giờ được mấy giờ $10$ phút?
  \shortans{6}
  \loigiai{
    $5\text{h}70 = 6\text{ giờ } 10\text{ phút}$.
  }
\end{ex}

% Câu 2
\begin{ex}
  Tính: $4\text{ ngày } 8\text{ giờ} - 1\text{ ngày } 14\text{ giờ}$. Kết quả có số giờ là bao nhiêu?
  \shortans{18}
  \loigiai{
    $3\text{ ngày } 32\text{ giờ} - 1\text{ ngày } 14\text{ giờ} = 2\text{ ngày } 18\text{ giờ}$.
  }
\end{ex}

% Câu 3
\begin{ex}
  Tính: $1\text{ giờ } 35\text{ phút} \times 3$. Đổi ra giờ được mấy giờ $45$ phút?
  \shortans{4}
  \loigiai{
    $3\text{ giờ } 105\text{ phút} = 4\text{ giờ } 45\text{ phút}$.
  }
\end{ex}

% Câu 4
\begin{ex}
  Tính: $8\text{ phút } 20\text{ giây} : 5$. Đổi ra giây là bao nhiêu giây?
  \shortans{100}
  \loigiai{
    $500\text{ giây} : 5 = 100\text{ giây}$.
  }
\end{ex}

% Câu 5
\begin{ex}
  Một người đi xe máy khởi hành lúc $7\text{ giờ } 15\text{ phút}$ và đến nơi lúc $9\text{ giờ } 45\text{ phút}$. Người đó đã đi hết bao nhiêu phút?
  \shortans{150}
  \loigiai{
    $2\text{ giờ } 30\text{ phút} = 150$ phút.
  }
\end{ex}

% Câu 6
\begin{ex}
  Một ô tô đi quãng đường $162\text{ km}$ hết $3$ giờ. Vận tốc của ô tô là bao nhiêu km/h?
  \shortans{54}
  \loigiai{
    $162 : 3 = 54\text{ km/h}$.
  }
\end{ex}

% Câu 7
\begin{ex}
  Một người đi xe đạp với vận tốc $15\text{ km/h}$ trong $2$ giờ $20$ phút ($2\frac{1}{3}$ giờ). Quãng đường đi được là bao nhiêu ki-lô-mét?
  \shortans{35}
  \loigiai{
    $15 \times \frac{7}{3} = 35\text{ km}$.
  }
\end{ex}

% Câu 8
\begin{ex}
  Hai xe ô tô cùng xuất phát từ hai đầu cầu dài $360\text{ km}$ đi ngược chiều nhau với vận tốc $50\text{ km/h}$ và $40\text{ km/h}$. Sau mấy giờ hai ô tô gặp nhau?
  \shortans{4}
  \loigiai{
    $360 : (50 + 40) = 4\text{ giờ}$.
  }
\end{ex}

% Câu 9
\begin{ex}
  Một người đi xe máy từ $A$ lúc $6\text{ giờ}$ với vận tốc $30\text{ km/h}$. Đến $7\text{ giờ}$, một ô tô từ $A$ đuổi theo với vận tốc $45\text{ km/h}$. Ô tô đuổi kịp xe máy lúc mấy giờ?
  \shortans{9}
  \loigiai{
    Khoảng cách lúc 7h: $30\text{ km}$. Hiệu vận tốc: $15\text{ km/h}$. Thời gian đuổi kịp: $30 : 15 = 2\text{ giờ}$. Lúc đó là $7 + 2 = 9\text{ giờ}$.
  }
\end{ex}

% Câu 10
\begin{ex}
  Một ca nô xuôi dòng với vận tốc $30\text{ km/h}$, ngược dòng với vận tốc $24\text{ km/h}$. Vận tốc của dòng nước là bao nhiêu km/h?
  \shortans{3}
  \loigiai{
    $(30 - 24) : 2 = 3\text{ km/h}$.
  }
\end{ex}

% Câu 11
\begin{ex}
  Ở bài ca nô trên, vận tốc thực của ca nô trong nước yên lặng là bao nhiêu km/h?
  \shortans{27}
  \loigiai{
    $(30 + 24) : 2 = 27\text{ km/h}$.
  }
\end{ex}

% Câu 12
\begin{ex}
  Quãng đường sông từ $A$ đến $B$ dài $48\text{ km}$. Ca nô đi xuôi dòng với vận tốc $24\text{ km/h}$. Thời gian ca nô đi hết bao nhiêu giờ?
  \shortans{2}
  \loigiai{
    $48 : 24 = 2\text{ giờ}$.
  }
\end{ex}

% Câu 13
\begin{ex}
  Một ô tô đi trong $3$ giờ đầu được $135\text{ km}$, trong $2$ giờ sau được $90\text{ km}$. Vận tốc trung bình của ô tô là bao nhiêu km/h?
  \shortans{45}
  \loigiai{
    $(135 + 90) : (3 + 2) = 225 : 5 = 45\text{ km/h}$.
  }
\end{ex}

% Câu 14
\begin{ex}
  Một người đi bộ $4$ vòng hồ hết $48$ phút. Muốn đi $6$ vòng hồ như thế thì hết bao nhiêu phút?
  \shortans{72}
  \loigiai{
    Mỗi vòng hết $48 : 4 = 12$ phút. $6$ vòng hết $12 \times 6 = 72$ phút.
  }
\end{ex}

% Câu 15
\begin{ex}
  Một con báo săn chạy với vận tốc $108\text{ km/h}$. Đổi vận tốc đó ra m/s được bao nhiêu m/s?
  \shortans{30}
  \loigiai{
    $108\,000 : 3\,600 = 30\text{ m/s}$.
  }
\end{ex}

% Câu 16
\begin{ex}
  Tính: $(2\text{ giờ } 40\text{ phút} + 1\text{ giờ } 20\text{ phút}) : 4$. Kết quả bằng đúng bao nhiêu giờ?
  \shortans{1}
  \loigiai{
    $4\text{ giờ} : 4 = 1\text{ giờ}$.
  }
\end{ex}

% Câu 17
\begin{ex}
  Một đoàn tàu đi qua một cây cột điện hết $12$ giây với vận tốc $15\text{ m/s}$. Chiều dài đoàn tàu là bao nhiêu mét?
  \shortans{180}
  \loigiai{
    $15 \times 12 = 180\text{ m}$.
  }
\end{ex}

% Câu 18
\begin{ex}
  Hai xe xuất phát cùng chiều: xe $A$ đi $60\text{ km/h}$, xe $B$ đi $40\text{ km/h}$. Sau $2$ giờ xe $A$ cách xa xe $B$ bao nhiêu ki-lô-mét?
  \shortans{40}
  \loigiai{
    $(60 - 40) \times 2 = 40\text{ km}$.
  }
\end{ex}

% Câu 19
\begin{ex}
  Bác Bình đi xe đạp $3\text{ km}$ hết $15$ phút. Vận tốc đi xe đạp của bác Bình là bao nhiêu km/h?
  \shortans{12}
  \loigiai{
    $15\text{ phút} = 0{,}25\text{ giờ}$. $3 : 0{,}25 = 12\text{ km/h}$.
  }
\end{ex}

% Câu 20
\begin{ex}
  Một xe buýt đi chặng đường dài $24\text{ km}$ với vận tốc $32\text{ km/h}$. Thời gian xe buýt chạy là bao nhiêu phút?
  \shortans{45}
  \loigiai{
    $24 : 32 = 0{,}75\text{ giờ} = 45$ phút.
  }
\end{ex}

% Câu 21
\begin{ex}
  Một vận động viên bơi $100\text{ m}$ hết $50$ giây. Vận tốc của vận động viên là bao nhiêu m/s?
  \shortans{2}
  \loigiai{
    $100 : 50 = 2\text{ m/s}$.
  }
\end{ex}

% Câu 22
\begin{ex}
  Một ô tô chở hàng chạy liên tục $4$ giờ được $180\text{ km}$. Nếu chạy $6$ giờ với cùng vận tốc thì được bao nhiêu ki-lô-mét?
  \shortans{270}
  \loigiai{
    Vận tốc: $180 : 4 = 45\text{ km/h}$. Sau 6 giờ: $45 \times 6 = 270\text{ km}$.
  }
\end{ex}

% Câu 23
\begin{ex}
  Một chuyến tàu hỏa xuất phát lúc $19\text{ giờ } 30\text{ phút}$ và đến nơi lúc $23\text{ giờ } 15\text{ phút}$. Thời gian chạy của tàu là bao nhiêu phút?
  \shortans{225}
  \loigiai{
    $23\text{h}15 - 19\text{h}30 = 3\text{ giờ } 45\text{ phút} = 225$ phút.
  }
\end{ex}

% Câu 24
\begin{ex}
  Hai người chạy bộ cùng chiều trên đường tròn dài $400\text{ m}$. Vận tốc người thứ nhất là $5\text{ m/s}$, người thứ hai là $3\text{ m/s}$. Sau bao nhiêu giây người thứ nhất hơn người thứ hai đúng $1$ vòng?
  \shortans{200}
  \loigiai{
    Hiệu vận tốc: $5 - 3 = 2\text{ m/s}$. Thời gian: $400 : 2 = 200\text{ giây}$.
  }
\end{ex}

% Câu 25
\begin{ex}
  Một người đi xe máy từ quê lên phố hết $2\text{ giờ } 15\text{ phút}$ với vận tốc $40\text{ km/h}$. Chiều dài quãng đường là bao nhiêu ki-lô-mét?
  \shortans{90}
  \loigiai{
    $40 \times 2{,}25 = 90\text{ km}$.
  }
\end{ex}
